\documentclass[10pt,a4paper]{amsart}
\usepackage[margin=19mm]{geometry}
\usepackage{amsmath,amssymb,mathtools}
\usepackage[scr=rsfs,cal=euler]{mathalfa}
\usepackage{xfrac}
\usepackage[unicode,hidelinks,bookmarks=false]{hyperref}
\newcommand{\ie}{i.e.\ }
\newcommand{\cf}{cf.\ }
\newcommand{\resp}{respectively\ }
\newcommand{\Subsection}{\subsection}
\providecommand{\nobreakdash}{\nobreak\hbox{-}}
\setlength{\emergencystretch}{2em}
\multlinegap0pt
\usepackage{bm}
\usepackage{bbm}
\usepackage{dsfont}
\usepackage{xspace}
\usepackage{cancel}
\usepackage{mathtools}
\usepackage{amsthm}
\makeatletter
\@ifundefined{theorem}{\newtheorem{theorem}{Theorem}}{}
\@ifundefined{lemma}{\newtheorem{lemma}[theorem]{Lemma}}{}
\makeatother
\DeclareMathOperator{\evacuation}{\mathcal{E}}
\title[Promotion permutations and the Robinson--Schensted correspon]{Promotion permutations and the Robinson--Schensted correspondence}
\author{Stephan Pfannerer, Joshua P. Swanson}
\date{Source: arXiv:2510.07744v1 (CC BY 4.0)}
\begin{document}
\maketitle

\noindent\emph{Source and licence.} Promotion permutations and the Robinson--Schensted correspondence. Version arXiv:2510.07744v1 (CC BY 4.0), distributed under the Creative Commons Attribution 4.0 International licence, as recorded in the arXiv licence metadata for this exact version and shown on the abstract page https://arxiv.org/abs/2510.07744v1. Full rights notice: licenses/arxiv-2510.07744-CC-BY-4.0.txt.\par\smallskip

\noindent\textit{Editorial scope.} Counted unit: the Lemma whose source label is \texttt{lem:proof-shadows} in the article (source0.tex lines 1120--1148, source label \texttt{lem:proof-shadows}), delivered together with the article's own complete original proof of it (source0.tex lines 1150--1209, one \texttt{proof} environment of 60 lines). No mathematics is written, repaired or re-derived: statement and proof are the article's own text line for line. Required context retained uncounted: none. Every definition, display and label of those lines is retained unchanged. Omitted from this package and not redistributed: 1--1119, 1149--1149, 1210--1222. Nothing inside a retained excerpt is omitted or altered: every retained body is delivered exactly as the article's lines are written, so no omission receipt of any kind accompanies this package. Citations: the retained excerpts carry no citation command at all, so no bibliography entry of the article is transcribed and no reference is redistributed. Reference adaptation: none was needed and none was made; every reference inside the delivered unit resolves inside it, so no unresolved reference or citation is produced. Printed slips kept literal: no printed word, symbol, hypothesis or number of the counted statement or of its proof was altered. Layout and class: the article's own class is \texttt{amsart} with packages including amsmath, amssymb, amsthm, thmtools, mathtools, array, enumitem, multicol, graphicx, xcolor, ytableau, geometry, microtype, cleveref; the wrapper is the lane's pinned \texttt{amsart} editorial wrapper at 10pt on A4, which restates the theorem-like environments the retained text needs and adds the article's own macro definitions that the retained text needs (\texttt{\textbackslash evacuation}), with extra wrapper packages (bm, bbm, dsfont, xspace, cancel, mathtools). The delivered numbering is the unit's own. Assets: no retained span contains a figure environment, a graphics-inclusion command, a PDF-page inclusion, a table or a TikZ picture; no image file is copied into this package, the package declares no asset dependency and no image file is redistributed with it. Licence: version arXiv:2510.07744v1 of this article is distributed under the Creative Commons Attribution 4.0 International licence, as recorded in the arXiv licence metadata for this exact version and shown on the abstract page https://arxiv.org/abs/2510.07744v1. A full rights notice is delivered at licenses/arxiv-2510.07744-CC-BY-4.0.txt. Characters: every retained excerpt is pure ASCII, so the delivered engine, XeLaTeX with its default Latin Modern text font, reports no missing character.
\begin{lemma}\label{lem:proof-shadows}
    The shadow lines of $S_k$ are
    \begin{align*}
    \begin{cases}
        \mathcal{L}_j(S_k, \nearrow) = S_{kj}
          & \text{where }1 \leq j \leq c\text{ and }r \leq k \leq 2r-1, \\
        \mathcal{L}_j(S_k, \swarrow) = S_{kj}
          & \text{where }1 \leq j \leq c\text{ and }1 \leq k \leq r.
    \end{cases}
    \end{align*}
    The shadow of $S_k$ is
    \begin{align*}
    \begin{cases}
        \mathcal{S}(S_k, \nearrow) = S_{k+1}
          & \text{where }r \leq k \leq 2r-1, \\
        \mathcal{S}(S_k, \swarrow) = S_{k-1}
          & \text{where }1 \leq k \leq r,
    \end{cases}
    \end{align*}
    where $S_0 = S_{2r} \coloneqq \varnothing$. Moreover,
    \begin{align*}
    \begin{cases}
        \mathcal{S}(S_{kj}, \nearrow) = S_{k+1,j}
          & \text{where }r \leq k \leq 2r-1, \\
        \mathcal{S}(S_{kj}, \swarrow) = S_{k-1,j}
          & \text{where }1 \leq k \leq r.
    \end{cases}
    \end{align*}
\end{lemma}

\begin{proof}
    We take $r \leq k \leq 2r-1$ throughout, the other case being analogous. Set $\ell \coloneq k-r+1 \leq r$. We have the following diagram of relations between the $(r-\ell+1) \cdot c$ entries of $S_k = S_{k1} \sqcup \cdots \sqcup S_{kc}$:
    \[
    \begin{array}{ccccccc}
        S_{k1}
          &
          &
          S_{k2}
          &
          &
          &
          & S_{kc} \\
        \hline
        (x_{\ell1},y_{r1})
          & \nearrow
          & (x_{\ell2},y_{r2})
          & \nearrow
          & \cdots
          & \nearrow
          & (x_{\ell c},y_{rc}) \\
        \searrow
          & \vdots
          & \searrow
          & \cdots
          & \searrow
          & \cdots
          & \searrow \\
        \vdots
          & \nearrow
          & \vdots
          & \nearrow
          & \ddots
          & \nearrow
          & \vdots\\
        \searrow
          & \vdots
          & \searrow
          & \cdots
          & \searrow
          & \cdots
          & \searrow \\
        (x_{r1},y_{\ell 1})
          & \nearrow
          & (x_{r2},y_{\ell 2})
          & \nearrow
          & \cdots
          & \nearrow
          & (x_{rc},y_{\ell c})
    \end{array}
    \]
    Here each $\nearrow$ and $\searrow$ relation follows directly from the row and column increasing conditions of $Q$ and $\evacuation(P)$.

    From these relations, we see directly that $\mathcal{L}_1(S_k, \nearrow) = S_{k1}$, and likewise for later shadows. The first claim follow. For the third, observe
      \[ (x_{\ell j}, y_{r j})
         \nearrow (x_{\ell+1,j}, y_{rj})
         \searrow (x_{\ell+1,j}, y_{r-1,j})
         \nearrow (x_{\ell+2,j}, y_{r-1,j})
         \searrow \cdots. \]
    This sequence interleaves $S_{kj}$ and $S_{k+1,j}$ in the required way. Now the second claim follows from the first and third.
\end{proof}

\end{document}
